\section{Introduction}
\label{sec:intro}

Unit tests are a practical safeguard against regressions, but writing them manually requires developers to identify representative inputs, boundaries, and expected outcomes. Code-oriented language models make direct test generation possible: a model can receive a function and produce a JUnit class in one call. The original HumanEval work established a widely used benchmark for code-generation evaluation \cite{chen2021codex}; however, test generation has an additional requirement. A generated test must be executable and its assertions must agree with the system under test (SUT). Text returned by an API is therefore only the first stage of the evaluation.

This study evaluates whether a fixed zero-shot prompt to GPT-4o-mini can generate useful JUnit~4 suites for Java functions. We use 63 functions from the repository's HumanEval-Java transformation. The study reports branch coverage and mutation score, rather than relying on a single coverage number. Branch coverage records whether decision outcomes were executed; mutation score records the proportion of injected changes killed by the suite. Together, the metrics distinguish broad execution from fault-detection strength.

We compare GPT-generated suites with retained EvoSuite suites. EvoSuite is a search-based Java test-generation tool \cite{fraser2011evosuite}; its existing 1-, 3-, and 5-minute archives were measured again using the same Maven, JaCoCo, and PIT pipeline. This is an operational technical comparison, not a student benchmark. The originally planned student-written comparator has no validated per-function measurements in the same pipeline and is therefore explicitly deferred.

The central research question is whether zero-shot GPT-4o-mini tests meet predefined coverage and mutation thresholds and how their metrics differ from EvoSuite on the same SUTs. We ask five subquestions: (RQ1) branch-coverage threshold attainment; (RQ2) mutation-score floor and target attainment; (RQ3) paired GPT--EvoSuite differences for each budget; (RQ4) simultaneous threshold attainment; and (RQ5) the technical failure pattern that blocks valid execution.

This paper contributes: (1) a traceable execution report for 63 Java targets, including raw API usage, repair records, and class-level metrics; (2) a paired, budget-specific comparison with existing EvoSuite suites; and (3) an explicit analysis boundary that prevents pass-conditioned results from being generalized to all generated suites. The repository retains the prompts, scripts, raw result files, and analysis notebook needed to reproduce the reported analysis.
